\section{Introduction}\label{sec:introduction}

Every infinite-dimensional Banach space contains a separable
infinite-dimensional closed subspace: take the closed span of a countable
linearly independent set. The corresponding assertion for quotients is the
\emph{separable quotient problem}, traditionally attributed to Banach and
Mazur. For a scalar field $\mathbb F\in\{\R,\C\}$, let $\SQ_{\mathbb F}$
denote the statement
\begin{quote}
Every infinite-dimensional Banach space over $\mathbb F$ admits a
bounded linear surjection onto a separable infinite-dimensional
Banach space.
\end{quote}
Equivalently, each such space $X$ has a closed subspace $M$ for which
$X/M$ is separable and infinite dimensional. The equivalence follows by
applying the open mapping theorem to the induced bijection from a space
modulo the kernel of a surjection.

We work with the axiom of Choice throughout. Write
$\mathfrak c=2^{\aleph_0}$ for the cardinality of the continuum.
For an uncountable cardinal $\kappa$, a probability measure on
$\mathcal P(\kappa)$ is \emph{$\kappa$-additive} if it is additive on
pairwise disjoint families of cardinality strictly less than $\kappa$;
the sum over such a family means the supremum of its finite subsums.
We say that $\kappa$ is \emph{real-valued measurable} if it carries such
a measure vanishing on singletons. In particular, the hypothesis below
requires a measure on \emph{every subset} of $\mathfrak c$, with this
stronger additivity property.

The two implications below are proved in ZFC. Their hypotheses are then
realized in models at the indicated consistency strengths; this gives
the relative-independence conclusion.

\begin{theorem}\label{intro:main}
For each $\mathbb F\in\{\R,\C\}$, ZFC proves the following implications.
\begin{enumerate}[label=\textup{(\roman*)}]
 \item If $\mathfrak c$ is real-valued measurable, then $\SQ_{\mathbb F}$
 holds.
 \item If the continuum hypothesis holds, then $\SQ_{\mathbb F}$ fails.
\end{enumerate}
Consequently, if ZFC with a measurable cardinal is consistent, then
$\SQ_{\mathbb F}$ is independent of ZFC. Consistency of its failure
requires only consistency of ZFC.
\end{theorem}

The CH counterexamples already occur at norm density $\aleph_1$, the
smallest possible nonseparable density (Corollary~\ref{twist:density}).
At this density, an earlier positive theorem yields a sharper consistency
conclusion. Let $\SQ^{\aleph_1}_{\mathbb F}$ denote the assertion that
every Banach space over $\mathbb F$ of norm density exactly $\aleph_1$
has a separable infinite-dimensional quotient.

\begin{corollary}\label{intro:small-density}
If ZFC is consistent, then $\SQ^{\aleph_1}_{\mathbb F}$ is independent
of ZFC for each $\mathbb F\in\{\R,\C\}$.
\end{corollary}

The positive side of this corollary uses the density criterion of Saxon
and S\'anchez Ruiz described below, together with Martin's axiom. The
negative side uses the counterexample constructed here. The measurable
cardinal hypothesis in Theorem~\ref{intro:main} concerns the unrestricted
quotient assertion.

\paragraph{History and methods.}
The history and many positive cases of the problem are surveyed
in~\cite{Brech2025,FerrandoKakolLopezPellicerSliwa2018}.
Brech notes that the problem is not stated explicitly in Banach's 1932
book and identifies Rosenthal's 1969 paper as an early written
formulation~\cite[p.~74]{Brech2025}.
Johnson and Rosenthal developed the weak-star basic-sequence method and
proved that every separable infinite-dimensional Banach space has an
infinite-dimensional quotient with a Schauder
basis~\cite{JohnsonRosenthal1972}.
Argyros, Dodos and Kanellopoulos proved the positive answer for every dual
Banach space~\cite[Theorem~15]{ArgyrosDodosKanellopoulos2008}.
Their Proposition~16, based on work of Hagler and
Johnson~\cite{HaglerJohnson1977}, also shows that an unconditional basic
sequence in $X^*$ gives a separable infinite-dimensional quotient of $X$
with an unconditional basis. This implication is central to our positive
argument. Morris and Yost give a common sufficient condition covering
both dual spaces and spaces that are the closed linear span of a weakly
compact subset~\cite{MorrisYost2020}.
More recently, Patri states a $c_0$-quotient criterion for
$\mathcal L_\infty$ spaces admitting an infinite-dimensional Schur
quotient, in which weakly convergent sequences converge in norm.
He applies the criterion to a class of nonseparable Bourgain--Pisier spaces~\cite{Patri2026}.

Earlier set-theoretic positive results depend on the density of the space.
Write $\mathfrak b$ for the least cardinality of a family in $\N^\N$
unbounded under eventual domination, where $f\leq^*g$ means
$f(n)\leq g(n)$ for all sufficiently large $n$.
Saxon and S\'anchez Ruiz proved that every infinite-dimensional Banach
space of density less than $\mathfrak b$ has a separable
infinite-dimensional quotient \cite{SaxonSanchezRuiz1996}; see also
\cite[Theorem~2.1]{Brech2025}. In particular, the positive assertion at
density $\aleph_1$ already follows under $\mathfrak b>\aleph_1$.
For large densities, Dodos, Lopez-Abad and
Todor\v{c}evi\'c proved, relative to infinitely many strongly compact
cardinals, consistency with GCH of a positive answer for every Banach
space of density at least $\aleph_\omega$
\cite{DodosLopezAbadTodorcevic2011}. The positive implication here applies
to every infinite-dimensional Banach space. The negative construction
uses conditional norming operations in the tradition of
Gowers--Maurey~\cite{GowersMaurey1993}, branch functionals as in
James-tree constructions~\cite{James1974,ArgyrosArvanitakisTolias2006},
and coding by order positions in coherent finite ordinal closures
\cite[Section~2.4]{ArgyrosLopezAbadTodorcevic2005}.
We give the particular construction and all its estimates in full.

We write $B_Z$ and $S_Z$ for the closed unit ball and unit sphere of a
Banach space $Z$, and $[z_i:i\in I]$ for a norm-closed linear span.
All separability and density assertions about subsets of dual spaces
refer to the norm topology unless another topology is specified.

\paragraph{The common criterion.}
For a closed subspace $F\subset X^*$, define the evaluation map
\[
 R_F:X\longrightarrow F^*,\qquad (R_Fx)(f)=f(x).
\]
A separable infinite-dimensional quotient exists precisely when some
infinite-dimensional closed $F$ has norm-separable evaluation image.
In the necessary direction $F$ may itself be chosen separable.
Section~\ref{sec:evaluation} proves this criterion, including the
surjectivity step.
The two constructions in the paper use this same criterion in opposite
ways.

For the positive implication, start with a putative counterexample $X$
and a separable infinite-dimensional subspace $E\subset X^*$.
The measure first forces $R_E(B_X)$ to have density $\mathfrak c$:
a smaller dense set would allow simultaneous summable evaluations on a
suitable basic sequence, and hence a separable quotient. From the large
evaluation image we construct a seminormalized family in $X^*$ which
vanishes almost everywhere at each point of $X$. Odell--Rosenthal supplies
sequential bidual approximations. Coding these approximations lets us use
ordinary Borel Fubini to select an unconditional basic sequence from that
family. This
contradicts the Hagler--Johnson consequence above. The Borel coding is
what permits the integration argument; no interchange theorem for
arbitrary power-set measures is needed.

For the negative implication, we build a dual space $E=X_0^*$ with
coordinate vectors indexed by $\omega_1$. We construct paths from
successive finitely supported functionals with bounded partial sums;
each path therefore defines a functional on $E$. Each piece carries
an integer weight. The coding makes distinct paths share only finitely
many weights, yielding uniform estimates for combinations of their
tails. A block argument
then gives, on every infinite-dimensional subspace of $E$, uncountably
many norm-separated restrictions of path functionals. CH lets us assign
all these paths to coordinates of $E$. A small perturbation of the
quotient map $E^*\to E^*/X_0$ produces a new predual $X$ whose evaluations
approximate all path cosets. Its evaluation image on every separable
infinite-dimensional subspace of $X^*\simeq E$ is nonseparable, excluding
all separable infinite-dimensional quotients. The perturbation lemma in
Section~\ref{sec:predual} isolates this change of predual from the
combinatorial construction.

\paragraph{Organization.}
Section~\ref{sec:evaluation} develops the evaluation criterion and a
selection lemma for strictly singular operators.
Section~\ref{sec:positive} proves the positive implication over $\R$,
and Section~\ref{sec:fields} transfers it to $\C$ while preparing the
corresponding transfer for counterexamples.
Sections~\ref{sec:norming}--\ref{sec:branches} construct the coordinate
dual, prove convergence for its block sequences, establish the required
block estimates, and separate functionals on arbitrary subspaces.
Section~\ref{sec:predual} changes its predual and proves the CH
counterexample. Section~\ref{model:section} supplies the set-theoretic models and
proves both Theorem~\ref{intro:main} and the sharper restricted-density
Corollary~\ref{intro:small-density}.
